\section{The transverse action and the fibre group}
\label{grp:section}

We use the intrinsic projection and the full incidence charts supplied by
Theorem~\ref{quo:projection}.  Write
\[
 p:V\longrightarrow X^+,\qquad f:V\longrightarrow M,
 \qquad \rho:\Gamma\longrightarrow\operatorname{Aut}(M),
 \qquad \Lambda=\rho(\Gamma).
\]
Here $p$ is a connected regular covering of a smooth projective manifold,
with group $\Gamma$, and $f$ is an equivariant surjective holomorphic
submersion with connected fibres.  As in the construction,
$\Gamma=\pi_1(X)$: the covering of $X^+$ is pulled back from the
universal covering of the original normal projective variety $X$.
The manifold $M$ is contractible and
biholomorphic to a bounded domain.  The complete real holomorphic vector
fields belonging to $\operatorname{Aut}(M)^0$ have evaluations spanning
$T_mM$ over $\mathbb C$, for every $m\in M$.

We recall precisely the compactification property used below.  Locally in
$M$, there is a smooth projective family $q:Q\to B$ with connected
fibres over a connected box $B$, a proper analytic
subset $D\subset Q$ contained in a relative divisor, and an open incidence
locus $G\subset Q$ with
\[
 Q\setminus D\subset G,
 \qquad e:G\longrightarrow f^{-1}(B),
 \qquad f\circ e=q|_G.
\]
After identifying the parameter box with its image in $M$, the map $e$
is surjective onto the whole of $f^{-1}(B)$ and is generically of full
rank.  The inclusion $Q\setminus D\subset G$ does not assert equality:
the points of $G\cap D$ will be retained.  Each individual fibre of $f$
also has a smooth projective compactification with analytic complement.
All assertions in this section depend on this full incidence property,
not only on a parametrization of a dense open subset of each leaf.

\subsection{The core and the external inputs}

A smooth compact K\"ahler manifold $Z$ is \emph{special} if it has no
Bogomolov sheaf: there is no line bundle $L$ with a nonzero morphism
$L\to\Omega_Z^r$ and $\kappa(Z,L)=r>0$.  We use the following two
consequences of Campana's core theory~\cite{Campana2004}.
There is a projective modification $r:Z'\to Z$ and a holomorphic core
$c:Z'\to C$ whose very general smooth connected fibre is special.
Its orbifold base is of general type unless $C$ is a point.  Moreover,
if $T$ is a connected complex manifold, $E$ is a reduced divisor, and
$a:T\setminus E\dashrightarrow Z$ is meromorphic, then
$c\circ r^{-1}\circ a$ extends meromorphically across $E$ whenever
this composition has maximal rank $\dim C$ somewhere.  This last
statement is the factorized form of the orbifold Kobayashi--Ochiai
theorem, namely \cite[Theorem~8.2]{Campana2004}.  It concerns the
orbifold core; the underlying variety $C$ need not itself be of general
type.  If $C$ is a point the extension assertion is immediate.

The other input is the compact special abelianity theorem of OpenAI,
\emph{The abelianity conjecture for special compact K\"ahler manifolds}
\cite[Theorem~1.1]{SpecialAbelianity}:
\begin{equation}\label{grp:abelianity-input}
 Z\text{ smooth, compact K\"ahler and special}
 \quad\Longrightarrow\quad \pi_1(Z)\text{ virtually abelian}.
\end{equation}
The conclusion concerns the entire ordinary fundamental group and has
no linearity or residual-finiteness hypothesis. We apply it to smooth
connected compact special core fibres, and to the smooth compact
manifold itself when its core is a point. All further conclusions
about the transverse action and its kernel are proved below.

\subsection{A Liouville lemma for abelian covers}

The following bounded-holomorphic conclusion is a consequence of
Lyons--Sullivan's theorem for nilpotent covers
\cite[Theorem~1, p.~300]{LyonsSullivan1984}, after passing to a finite
cover of the compact base. We retain a self-contained proof for the
abelian case needed here.

\begin{lemma}\label{grp:liouville}
Let $a:W\to Z$ be a connected regular covering of a compact connected
K\"ahler manifold.  If its deck group is virtually abelian, every bounded
holomorphic function on $W$ is constant.
\end{lemma}

\begin{proof}
Choose an abelian subgroup $A$ of finite index in the deck group.
The quotient $Z_A=W/A$ is a compact K\"ahler manifold.  Pull a K\"ahler
metric on $Z_A$ back to $W$, and let $\Delta=\operatorname{div}\nabla$.
The pullback metric is complete.  Fix $o\in W$ and consider the convex set
\[
 \mathcal H=\{h\in C^\infty(W):h>0,\ \Delta h=0,\ h(o)=1\}.
\]
Local Harnack inequalities
\cite[Theorems~5--7]{Serrin1964} and interior elliptic estimates
\cite[Lecture~IV, equation~(4.2)$'$ and its consequences]{Nirenberg1959}
make $\mathcal H$ compact in the topology of smooth convergence on
compact subsets. In relatively compact metric charts the scalar
Laplacian has smooth uniformly elliptic coefficients; its lower-order
terms in divergence form vanish, so these estimates apply without
an additive Harnack term.
Limits are still strictly positive: they are nonnegative harmonic
functions equal to one at $o$, so the strong maximum principle applies.

Let $h$ be an extreme point of $\mathcal H$.  Its ray is minimal in the
cone of positive harmonic functions.  Indeed, if $0\leq v\leq h$ is
harmonic and $0<v(o)<1$, then
\[
 h=v(o)\frac{v}{v(o)}+(1-v(o))\frac{h-v}{1-v(o)},
\]
so extremality makes $v$ a multiple of $h$.  The endpoint cases follow
from the maximum principle.

For each $\alpha\in A$ there is a number $C_\alpha$ such that
\begin{equation}\label{grp:harnack-translate}
 h(\alpha x)\leq C_\alpha h(x)\qquad(x\in W)
\end{equation}
for every positive harmonic function $h$.  To see the uniformity, choose
a compact set $K\subset W$ with $AK=W$.  Harnack chains between the two
compact sets $K$ and $\alpha K$ give one constant for all positive
harmonic functions.  If $x=\beta k$, where $\beta\in A$ and $k\in K$,
apply this inequality to $h\circ\beta$ and use
$\alpha\beta=\beta\alpha$.  This proves
\eqref{grp:harnack-translate}.  Minimality of the ray of $h$ now gives
$h\circ\alpha=c_\alpha h$ for some $c_\alpha>0$.

Consequently $\eta=d\log h$ is $A$-invariant and descends to a smooth
one-form on $Z_A$.  Its metric dual satisfies
\[
 \operatorname{div}(\eta^\sharp)
 =\Delta\log h=-|\eta|^2.
\]
Integrating on $Z_A$ shows that $\eta=0$.  Notice that the differential
descends; a single-valued function $\log h$ on $Z_A$ is not required.
Thus every extreme point of $\mathcal H$ is the constant function one.
The Krein--Milman theorem implies $\mathcal H=\{1\}$.
Finally, the real and imaginary parts of a holomorphic function are
harmonic for a K\"ahler metric.  Adding a constant to each bounded real
part makes it positive, so both parts are constant.
\end{proof}

\subsection{From countable levels to locally finite parameters}

The following elementary analytic observation records the exact role of
properness.  Countability by itself would not give discreteness.

\begin{lemma}\label{grp:finite-level-projection}
Let $q:Q\to B$ be a proper holomorphic map from a complex manifold, let
$h:Q\dashrightarrow C$ be a meromorphic map to a projective variety,
and let $G_0\subset Q$ be open with $h|_{G_0}$ holomorphic.  Fix $s\in C$.
If
\[
 q\bigl(G_0\cap h^{-1}(s)\bigr)
\]
is countable, it is locally finite in $B$.
\end{lemma}

\begin{proof}
Resolve the graph of $h$, obtaining a proper modification
$\tau:\widehat Q\to Q$ and a holomorphic map
$\widehat h:\widehat Q\to C$.
Put $\widehat q=q\tau$ and $A=\widehat h^{-1}(s)$, with its reduced
analytic structure.  Over $G_0$, equality of meromorphic maps gives
$\widehat h=h\tau$.  In particular every point of
$G_0\cap h^{-1}(s)$ has preimages in
$A\cap\tau^{-1}(G_0)$.

Let $A_i$ be an irreducible component of $A$ meeting
$\tau^{-1}(G_0)$.  On its nonempty relatively open part in this set,
$\widehat q$ has countable image.  A holomorphic map of positive rank
on a smooth patch has uncountable image by the holomorphic rank theorem.
Thus $\widehat q$ has rank zero on the regular locus of that open part
and is constant on some nonempty relatively open subset of $A_i$.
The identity theorem on an irreducible reduced complex space then makes
$\widehat q$ constant on all of $A_i$.  Equivalently, a fibre of
$\widehat q|_{A_i}$ contains an open subset and hence equals $A_i$.

For a relatively compact smaller box $B'\Subset B$,
$\widehat q^{-1}(\overline{B'})$ is compact.  The irreducible components
of a complex analytic space are locally finite, so only finitely many
$A_i$ meet this compact set.  Every attained parameter in $B'$ is the
constant value of one of the components meeting $\tau^{-1}(G_0)$.
There are therefore only finitely many such parameters.  Components
lying entirely outside the valid open set contribute no parameter and
need not be excluded from the compact component count.
\end{proof}

Apply the core construction to $X^+$.  Choose a smooth projective
modification $r:Z\to X^+$ and a holomorphic core $c:Z\to C$.  There is
a Zariski open subset $X^\circ\subset X^+$ over which $r$ is an
isomorphism, the induced core map is holomorphic, and its restriction
has full rank.  Shrink this open set if necessary so that very general
fibres meet it.  For a very general $s\in C$ put
\[
 F_s=c^{-1}(s),\qquad
 F_s^\circ=F_s\cap r^{-1}(X^\circ),\qquad
 D_s=f\bigl(p^{-1}(r(F_s^\circ))\bigr)\subset M.
\]
The variety $F_s$ is smooth, connected, projective and special.

\begin{lemma}\label{grp:countable-core-level}
For every such $s$, the set $D_s$ is countable and locally finite in $M$.
The union of these sets, as $s$ runs through very general values, is
dense in $M$.
\end{lemma}

\begin{proof}
Pull the regular covering $p$ back along $F_s\to X^+$.
Each connected component of the pullback is a regular cover of $F_s$
with deck group the image of $\pi_1(F_s)$ in $\Gamma$, up to the
basepoint identification.  That image is virtually abelian by
\eqref{grp:abelianity-input}.  The composite of the lifted map to $V$
with each bounded coordinate of $f$ is constant by
Lemma~\ref{grp:liouville}.  Hence $f$ is constant on each connected
component of the pullback.  There are at most countably many components,
since the fibre of $p$ is countable.  This proves countability of $D_s$.

Fix a full incidence chart $q:Q\to B$, $e:G\to f^{-1}(B)$.
On $Q\setminus D$ consider
\[
 a=p\circ e:Q\setminus D\longrightarrow X^+.
\]
Replace the containing analytic subset by a reduced divisor if needed.
The total space $Q$ is connected.  Its full-rank evaluation locus is a
nonempty open set; removing the containing divisor leaves a nonempty
open subset on which $a$ is submersive.  Its image contains an open
subset of $X^+$ and thus meets the dense ordinary core locus $X^\circ$.
At a preimage of such a point, the composition with the core has rank
$\dim C$.  The extension theorem therefore gives a
meromorphic map
\[
 h:Q\dashrightarrow C,
 \qquad h|_{Q\setminus D}=c\circ r^{-1}\circ p\circ e.
\]
On the \emph{full} open incidence locus $G$, uniqueness of meromorphic
continuation identifies this extension with the same meromorphic
composition.  In particular it is holomorphic and agrees with that
composition on
\[
 G_0=e^{-1}\bigl(p^{-1}(X^\circ)\bigr)\subset G.
\]
This identity also holds at points of $G_0\cap D$: they have not been
removed from the argument.

Surjectivity of $e$ onto entire leaves gives the exact equality
\begin{equation}\label{grp:full-incidence-level}
 D_s\cap B=q\bigl(G_0\cap h^{-1}(s)\bigr).
\end{equation}
Indeed a point on either side has a representative in the full incidence
locus whose image in $X^\circ$ belongs to the ordinary core fibre at
$s$.  Lemma~\ref{grp:finite-level-projection}, applied to
\eqref{grp:full-incidence-level}, proves local finiteness.  Its proof
does not require that additional exceptional sets in $G_0$ extend to
the compactification: $G_0$ is an open set, and every compactified level
component meeting it already has constant parameter.  Nor does it
require evaluation to be submersive at every representative of an
attained parameter.

Very general values omit a countable union of proper analytic subsets
of $C$.  On the ordinary submersion locus their inverse images have
dense complement, by the Baire theorem in coordinate boxes.  The
ordinary locus is dense in $X^+$; the covering $p$ and the submersion
$f$ are open.  Their images therefore show that the union of the $D_s$
is dense in $M$.  When $C$ is a point use the unique level; the
countability argument and connectedness already force $M$ to be a point.
\end{proof}

\begin{proposition}\label{grp:discrete}
The image $\Lambda$ is discrete in $\operatorname{Aut}(M)$ and acts
cocompactly on $M$.
\end{proposition}

\begin{proof}
The group of automorphisms of a bounded domain is a real Lie group
\cite[Proposition~3.1 and Theorem~6.2]{Kobayashi1967} and acts properly
\cite[\S2, Propositions~2--3]{BlanchardCartan1954}.  Let $H$ be the closure of $\Lambda$ in this group.
For every very general $s$, the set $D_s$ is $\Lambda$-invariant: deck
transformations preserve $p$, and $f$ is equivariant.  A locally finite
subset of a manifold is closed, so $D_s$ is also $H$-invariant by
continuity.  The orbit under the connected group $H^0$ of any point of
$D_s$ is connected and lies in a discrete set.  It is a point.
Thus $H^0$ fixes the dense union of the $D_s$ pointwise and consequently
fixes $M$ pointwise.  It is the trivial subgroup.  As a closed subgroup
of a Lie group, $H$ is a Lie group
\cite[Chapter~II, \S III, no.~27]{CartanLie1952}; zero-dimensional Lie groups are
discrete.  This proves the assertion for $\Lambda$.

The equivariant surjection $f$ induces a continuous surjection
$X^+=V/\Gamma\to M/\Lambda$.  Its source is compact, so the quotient
is compact.
\end{proof}

\subsection{Removing stabilizers and identifying the domain}

\begin{lemma}\label{grp:free-finite-index}
There is a torsion-free finite-index subgroup $\Lambda_1\subset\Lambda$
acting freely on $M$.  Its inverse image in $\Gamma$ corresponds to a connected
finite cover of the original normal projective variety $X$.
\end{lemma}

\begin{proof}
The group $\Gamma$ is finitely generated, being a quotient of
$\pi_1(X^+)$, and hence so is $\Lambda$.
Let $\mathfrak g$ be the real Lie algebra of $\operatorname{Aut}(M)$.
Selberg's lemma \cite[Lemma~8]{Selberg1960}, in the matrix-group
form recalled in \cite[p.~435]{Wolf1968}, applied to the finitely
generated linear group
$\operatorname{Ad}(\Lambda)\subset\operatorname{GL}(\mathfrak g)$
gives a torsion-free subgroup $L$ of finite index.  Set
\[
 \Lambda_1=(\operatorname{Ad}|_\Lambda)^{-1}(L).
\]
For $m\in M$, properness and discreteness make $(\Lambda_1)_m$ finite.
Its image in the torsion-free group $L$ is trivial.
If $\lambda\in(\Lambda_1)_m$, therefore, $\lambda_*\xi=\xi$ for every
complete vector field $\xi\in\mathfrak g$.  Evaluating at the fixed
point gives
\[
 d\lambda_m\bigl(\xi(m)\bigr)=\xi(m).
\]
The complex spanning property and complex linearity of the derivative
give $d\lambda_m=1$.  A finite-order holomorphic automorphism fixing a
point with identity derivative is the identity: if its order is $k$
and $z$ is a local coordinate system centered there, then
$k^{-1}\sum_{j=0}^{k-1}z\circ\lambda^j$ is an invariant coordinate
system, since its derivative is the identity.  The inverse function
theorem and analytic continuation prove $\lambda=1$ on $M$.
Thus all stabilizers are trivial.  No faithfulness of the full adjoint
representation has been assumed.

The subgroup $\Lambda_1$ is also torsion-free.  Indeed, lift a finite
triangulation of the compact smooth quotient $M/\Lambda_1$ to its
contractible universal cover.  Its cellular chains form a finite-length
free resolution of $\mathbb Z$ over $\mathbb Z[\Lambda_1]$.  Restriction
to a subgroup of prime order would remain a free resolution of finite
length, contrary to the nonzero cohomology in arbitrarily high even
degrees of a finite cyclic group.  Every nontrivial finite-order element
has a power of prime order, so no such element exists.

The subgroup $\rho^{-1}(\Lambda_1)$ has finite index in $\Gamma$.
The corresponding connected topological covering of $X$ is a finite
unramified analytic covering. The proper-base Riemann existence
theorem \cite[XII, Corollary~4.6]{SGA1} algebraizes it. Pullback
of an ample bundle under the resulting finite morphism is ample
\cite[Tag~0B5V]{StacksProject}, so the cover is projective;
its analytic covering charts make it normal.  This justifies the simultaneous finite-cover replacement.
\end{proof}

\begin{theorem}\label{grp:discrete-symmetric}
Under the intrinsic-projection hypotheses recalled above, the transverse
image is discrete and cocompact.  After a finite cover it acts freely,
and $M$ is biholomorphic to a bounded symmetric domain, with a point
allowed.
\end{theorem}

\begin{proof}
Discreteness and the free finite-index action were just proved.
Replace $\Lambda$ by the subgroup supplied by
Lemma~\ref{grp:free-finite-index} and put $B=M/\Lambda$.
The intrinsic Bergman kernel of square-integrable canonical forms on
the bounded realization of $M$ is everywhere positive, and its logarithm
has positive definite complex Hessian.  The reciprocal kernel is a
Hermitian metric on $K_M$; invariance under biholomorphisms makes it
descend to a metric on $K_B$ with positive curvature.  Consequently
$K_B$ is ample by the Kodaira embedding theorem
\cite[\S\S1--2; see also \S4]{Kodaira1954PNAS}, $B$ is projective,
and $c_1(B)<0$. In complex dimension one the same positivity-to-ampleness
step follows from Riemann--Roch; the point case is immediate.

The manifold $B$ is compact and aspherical because its universal cover
$M$ is contractible.  The Nadel--Frankel splitting theorem in the
aspherical form of \cite[Theorem~1.10]{FarbWeinberger2008} applies.
It gives, on universal covers, a holomorphic decomposition
\[
 M\simeq D_{\mathrm{sym}}\times R,
\]
where $D_{\mathrm{sym}}$ is a bounded symmetric domain and
$\operatorname{Aut}(R)$ is discrete.  The cited statement is obtained
from a product decomposition of a finite cover of $B$; this does not
change its universal cover.  Its Hermitian locally symmetric factor is
of noncompact type, as its canonical bundle is ample, so its simply
connected cover has the bounded symmetric realization used here.

We explain why the identity automorphisms of $M$ cannot mix these two
factors.  The Bergman metric of $M$ descends to compact $B$, so is
complete.  On a product, Fubini's theorem identifies the Hilbert space
of square-integrable canonical forms with the Hilbert tensor product
of the two factor spaces.  Taking product orthonormal bases gives the
product formula for the kernels.  It follows that
\[
 g_M=g_{D_{\mathrm{sym}}}\oplus g_R.
\]
Both factor metrics are positive definite, since $g_M$ is, and complete,
since a factor slice is a closed totally geodesic submanifold of this
complete product.  Both factors are simply connected, since $M$ is.

There is no Euclidean de Rham factor in $M$.  Otherwise its parallel
vector fields span a nonzero parallel real subspace preserved by the
parallel complex structure, hence yield a holomorphic Euclidean factor
$\mathbb C^a$, $a>0$.  Restricting a bounded realization of $M$ to a
complex line in this factor contradicts Liouville's theorem.  Uniqueness
of the non-Euclidean irreducible de Rham factors
\cite[\S7, Th\'eor\`eme~III]{DeRham1952} now implies that every
isometry in the identity component preserves each factor distribution,
and in particular the two displayed factor foliations.  Thus every
element of $\operatorname{Aut}(M)^0$ acts as a product automorphism.
Its component on $R$ is constant as a function of the group element,
because $\operatorname{Aut}(R)$ is discrete; it is the identity.
Every vector field from this identity group therefore has zero
$R$-component.  The complex spanning property forces $\dim R=0$.
Hence $M\simeq D_{\mathrm{sym}}$.
\end{proof}

\subsection{Compactifiable covers and finite-index fibre images}

\begin{lemma}\label{grp:compactifiable-cover}
Let $a:U\to N$ be a connected regular covering of a smooth connected
projective manifold, with deck group $K$.  Suppose $U$ is a Zariski open
subset, in the analytic sense, of a smooth projective manifold $\overline U$.
Then, using \eqref{grp:abelianity-input}, the group $K$ is virtually
abelian.
\end{lemma}

\begin{proof}
Choose a smooth projective core model, with $C$ smooth,
\[
 r:N'\longrightarrow N,\qquad c:N'\longrightarrow C.
\]
Let $N^\circ\subset N$ be a Zariski open set on which $r$ is an
isomorphism and the ordinary core is holomorphic and submersive.
The composition $c\circ r^{-1}\circ a$ extends meromorphically to
$\overline U$ by the core extension theorem: the boundary is analytic,
it can be included in a reduced divisor, and $a$ is locally an
isomorphism.  Resolve this meromorphic map to obtain
\[
 \tau:\widehat U\longrightarrow\overline U,
 \qquad h:\widehat U\longrightarrow C.
\]
The modification can be chosen to be an isomorphism above
$a^{-1}(N^\circ)$.  Put $\widehat U_U=\tau^{-1}(U)$.
The complement of $\widehat U_U$ is an analytic subset of the compact
smooth projective manifold $\widehat U$.

Take $s\in C$ very general, also outside the proper analytic sets
required for generic smoothness of $h$ and $c$.  The compact level
$L=h^{-1}(s)$ has finitely many connected components, each a smooth
compact complex manifold.  For every component $L_i$ meeting
$\widehat U_U$, its complement of the analytic boundary is connected.
Within this connected complex manifold, the further excluded set
\[
 (a\circ\tau)^{-1}(N\setminus N^\circ)
\]
is analytic.  If the component meets the retained ordinary locus,
this is a proper analytic subset and its removal is still connected.
Here we use the elementary fact that a proper analytic subset of a
connected complex manifold does not disconnect it: a path can be
perturbed to avoid its real-codimension-at-least-two stratification.
The excluded subset need not extend to all of $L_i$ for this argument.
Components contained in the excluded set are simply discarded.

It follows that the retained ordinary level
\begin{equation}\label{grp:finite-ordinary-level}
 L\cap\tau^{-1}(a^{-1}(N^\circ))
 \simeq a^{-1}(r(F_s^\circ)),
 \qquad
 F_s=c^{-1}(s),\quad F_s^\circ=F_s\cap r^{-1}(N^\circ),
\end{equation}
has finitely many connected components.  The isomorphism uses the
choice that $\tau$ is an isomorphism on the ordinary locus.  The fibre
$F_s$ is smooth, connected, projective and special, and $F_s^\circ$ is
the complement of a proper analytic subset of it.

Fix a point of $F_s^\circ$ and a lift to $U$.  The monodromy
representation of the regular covering $a$ identifies its fibre with
$K$.  If
\[
 H=\operatorname{im}\bigl(\pi_1(F_s^\circ)\longrightarrow K\bigr),
\]
the connected components of the restricted covering in
\eqref{grp:finite-ordinary-level} are the $H$-orbits in this fibre,
equivalently the cosets of $H$ in $K$.  Thus
\begin{equation}\label{grp:index-equals-components}
 [K:H]=\#\pi_0\bigl(a^{-1}(r(F_s^\circ))\bigr)<\infty.
\end{equation}

There is no increase of the relevant monodromy when the compact fibre
is replaced by $F_s^\circ$.  Indeed the map to $N$ extends as
$r|_{F_s}:F_s\to N$.  Moreover
$\pi_1(F_s^\circ)\to\pi_1(F_s)$ is surjective, by perturbing loops
off the proper analytic complement.  Hence
\[
 H=\operatorname{im}\bigl(\pi_1(F_s)\longrightarrow\pi_1(N)
                  \longrightarrow K\bigr).
\]
The group $\pi_1(F_s)$ is virtually abelian by
\eqref{grp:abelianity-input}; so is its quotient $H$.
An abelian subgroup of finite index in $H$ also has finite index in
$K$, by \eqref{grp:index-equals-components}.  This proves the lemma.

If the core is a point, one may instead apply
\eqref{grp:abelianity-input} directly to $N$, since $N$ is then special
and $K$ is a quotient of $\pi_1(N)$.  This also covers dimension zero.
\end{proof}

\begin{theorem}\label{grp:kernel-abelian}
After the finite-cover replacement of
Lemma~\ref{grp:free-finite-index}, let
$K=\ker(\Gamma\to\Lambda)$.  The group $K$ is finitely generated and
virtually abelian.  For each $m\in M$, it acts on $f^{-1}(m)$ as the
deck group of a connected regular covering of a smooth projective
fibre.
\end{theorem}

\begin{proof}
Since $\Lambda$ acts freely, the stabilizer in $\Gamma$ of $m$ is
exactly $K$.  The equivariant submersion descends to a holomorphic
submersion
\[
 g:X^+\longrightarrow B=M/\Lambda.
\]
It is proper because $X^+$ is compact.  Its fibre $N_b$, for the image
$b$ of $m$, is a smooth connected projective manifold.  Restricting the
covering $p$ gives
\[
 f^{-1}(m)\longrightarrow N_b
\]
as a connected regular covering with group $K$; all identifications
follow by lifting to a small evenly covered neighborhood in $B$.
The fibre compactifications in Theorem~\ref{quo:projection} give
the hypothesis of Lemma~\ref{grp:compactifiable-cover}.  If the initial
compactifying fibre is singular, it is smooth on the open leaf;
resolve its boundary while leaving that leaf unchanged.  The lemma
therefore proves that $K$ is virtually abelian.
Finally the monodromy map $\pi_1(N_b)\to K$ is surjective, because the
restricted covering is connected and regular.  A compact manifold has
finitely generated fundamental group, so $K$ is finitely generated.
\end{proof}
